\documentclass[exercice]{thedocument}%
\usepackage{tabularray}%
\startdatakeys
\source{}
\BO{2026}
\cycle{4}
\competencesDuSocle{RE,RA,MO}
\objectifsApprentissage{C4D1bb,C4D1bc,C4D2ah}
\automatismes{C4C2ab,C4D3ab}
\enddatakeys
\begin{document}%
Chez le marchand de fruits, on voit le tableau suivant.\vskip10pt
\begin{tblr}{%
        colspec={Q[m,wd=80pt]*3{Q[c,m,wd=30pt]}},
        column{1}={font=\bfseries,bg=gray!25!white},
        hlines,
        vlines}
    Masse de pommes {\normalfont( en kg)}&2&4&5\\
    Prix {\normalfont( en euros)}&5,00&9,00&10,00
\end{tblr}\vskip10pt\noindent
Le prix des pommes est-il proportionnel à la masse achetée~? Justifier la réponse.
\end{document}
